% Generated draft: requires mathematical and source review.
\documentclass[11pt]{article}
\usepackage[margin=25mm]{geometry}
\usepackage{fontspec}
\setmainfont{Latin Modern Roman}
\usepackage{amsmath,amssymb,mathtools,unicode-math}
\setmathfont{Latin Modern Math}
\usepackage{enumitem,needspace,xurl,hyperref}
\setlength{\parindent}{0pt}
\setlength{\parskip}{5pt}
\setlength{\emergencystretch}{3em}
\newcommand{\N}{\mathbb{N}}
\newcommand{\Z}{\mathbb{Z}}
\newcommand{\Q}{\mathbb{Q}}
\newcommand{\R}{\mathbb{R}}
\newcommand{\C}{\mathbb{C}}
\newcommand{\F}{\mathbb{F}}
\newcommand{\A}{\mathbb{A}}
\newcommand{\PP}{\mathbb{P}}
\newcommand{\E}{\mathbb{E}}
\newcommand{\eps}{\varepsilon}
\newcommand{\dd}{\,\mathrm{d}}
\begin{document}
\section*{UnsolvedMath ID: 30001879}
\textbf{Source problem code:} catalogue entry
\subsection{Definitions and mathematical statement}
Let \(S\subset\R^3\) be a finite set of \(n\) distinct points. A line determined by \(S\) is an affine line containing at least two points of \(S\); a plane determined by \(S\) is an affine plane containing at least three noncollinear points of \(S\). Define
\[
 l(S)=\left|\{\operatorname{aff}\{x,y\}:x,y\in S,\ x\ne y\}\right|,
\]
\[
 p(S)=\left|\{\operatorname{aff}\{x,y,z\}:x,y,z\in S,\ x,y,z\text{ are noncollinear}\}\right|.
\]
Here \(\operatorname{aff}\) denotes affine span, and distinct lines or planes are counted once. Two affine lines are skew if they are not contained in a common affine plane.

\textbf{Conjectural statement attributed to Purdy in the supplied record.} Does there exist an integer \(n_0\ge 3\) such that, for every integer \(n\ge n_0\) and every set \(S\subset\R^3\) of cardinality \(n\), the hypotheses
\[
 S\not\subset L_1\cup L_2\quad\text{for every pair of skew affine lines }L_1,L_2,
\]
and
\[
 |S\cap H|\le n-2\quad\text{for every affine plane }H\subset\R^3
\]
imply
\[
 p(S)\ge l(S)?
\]
The plane hypothesis interprets ``no plane containing \(n-1\) points'' as forbidding a plane containing at least \(n-1\) points. If the source intended to forbid only planes containing exactly \(n-1\) points, this formulation would differ and requires correction after source review.

\textbf{Unresolved source qualification.} The supplied statement fields give the question above, and the supplied verification metadata says that it was repaired against the official report. However, the supplied background and literature assessment say that only the historical attribution of Conjecture P3 to Purdy was extracted and that its full target and hypotheses are absent. This conflict is unresolved here. The displayed formulation is a statement-only reconstruction of the supplied question, requiring human mathematical and source review.
\subsection{Short English statement}
The supplied question asks whether every sufficiently large finite point set in \(\R^3\), subject to the stated exclusions involving two skew lines and a plane through all but at most one point, determines at least as many planes as lines. The record attributes the conjecture to Purdy and labels it open, but its conflicting source qualifications require review.
\subsection{Sources}
\begin{itemize}
\item Local registry: \texttt{data/unsolvedmath/problems.json}, permanent UnsolvedMath ID \(30001879\), source problem code \texttt{catalogue entry}, title ``Purdy Plane Line Incidence Conjecture.'' The supplied record belongs to the Oberwolfach Reports Open Problems collection and records proposed year \(2011\); its proposer field is empty. Attribution to Purdy and the designation Conjecture P3 occur in the supplied background and literature assessment.
\item Supplied report citation: \emph{Discrete Geometry (2012)}, Oberwolfach Reports, pp.\ 2459--2548, DOI \texttt{10.4171/owr/2011/44}. Supplied source URL: \url{https://doi.org/10.4171/owr/2011/44}. The literature-sources entry labels this same link ``original OWR report.'' The supplied title year \(2012\) and the \(2011\) year appearing elsewhere in the record are retained without reconciliation.
\item Supplied literature assessment dated \texttt{2026-08-21}: the status is labeled ``open,'' while the assessment says that only the historical statement that Purdy formulated Conjecture P3 was extracted, that the full target and hypotheses are absent, and that the problem therefore cannot be certified solved or disproved. This is supplied dated curation, not an independently verified literature assessment. Its instruction to verify time-sensitive status and cited claims independently is retained as a review requirement.
\item Supplied statement-verification metadata labels the statement \texttt{corrected\_verified} and reports that the official report was checked and the statement repaired to match its passage. These are claims of the supplied record; no independent source inspection or verification is asserted here.
\end{itemize}
\end{document}
